% SPDX-FileCopyrightText: 2026 Hasan Melih Akbulut
% SPDX-License-Identifier: CC-BY-4.0

\chapter{The measurements, together}
\label{ch:results}

The preceding chapters each carried their own numbers. This one puts
them on one page, states what each of them is evidence \emph{of}, and
then reports the four experiments that were run while this document was
being written, because they are the newest measurements in the project
and two of them change what an earlier chapter would otherwise have
said.

\section{The state of the part, on one page}

\Cref{tab:headline} is the whole design in numbers. Every row names the
document that owns the measurement; where a figure has been superseded,
the superseding document is named rather than the figure silently
replaced, which is the rule \dcite{64}{} sets for this project.

\begin{table}[tbp]
\centering
\caption{Headline measurements of the design as of 15 September 2026.
Each row is measured, not estimated, and names its source. The
layout rows are all from the sign-off run \code{s83romecc5}, the only
layout in the project that meets hold at every corner.}
\label{tab:headline}
\small
\begin{tabularx}{\linewidth}{@{}lXr@{}}
\toprule
\textbf{What} & \textbf{Value} & \textbf{Source} \\
\midrule
\multicolumn{3}{@{}l}{\textit{Process and floorplan}}\\
Technology & IHP SG13G2, 130\,nm, open PDK & \code{docs/12} \\
Die & $2306.40 \times 2075.22$\,\si{\micro\metre} & final DEF \\
Core area & $(60.00,\,45.36)$ to $(2246.40,\,2029.86)$\,\si{\micro\metre} & \code{clint\_region.py} \\
Vendor macros & 8: four $2048\times64$ RAM, two $1024\times32$ ROM, two $512\times16$ ROM ECC & \code{docs/79} \\
Instances placed & 168{,}592 & \code{metrics.json} \\
Instance utilisation & 0.7538 (standard cells 0.4662) & \code{metrics.json} \\
Flip-flops & 5{,}998 & final netlist \\
\midrule
\multicolumn{3}{@{}l}{\textit{Timing, three corners, no frequency claimed}}\\
Worst hold slack, \code{nom\_fast\_1p32V\_m40C} & $+0.0296$\,ns & \code{metrics.json} \\
Worst hold slack, \code{nom\_typ\_1p20V\_25C} & $+0.1769$\,ns & \code{metrics.json} \\
Worst hold slack, \code{nom\_slow\_1p08V\_125C} & $+0.4261$\,ns & \code{metrics.json} \\
Worst setup slack, slow corner & $-7.5758$\,ns & \code{metrics.json} \\
Setup TNS, slow corner & $-14{,}735.65$\,ns over 3{,}529 endpoints & \code{violator\_census.py} \\
Dominant violating launch point & the register file's read address & \code{docs/83} \\
Dominant violating capture class & the register file and the rest of Ibex, 2{,}088 of 3{,}529 endpoints & \code{docs/83} \\
\midrule
\multicolumn{3}{@{}l}{\textit{Geometric sign-off}}\\
Route DRC errors & 0 & \code{metrics.json} \\
Antenna violating nets & 2, and clearing them breaks hold & \code{docs/79} \\
LVS & matches uniquely, macros as black boxes & \code{docs/67} \\
XOR & 0 differing shapes & \code{docs/67} \\
\midrule
\multicolumn{3}{@{}l}{\textit{Where the logic sits} (\cref{fig:layout_panels})}\\
Ibex core and its register file & 23{,}942 instances, upper channel & this document \\
Neuromorphic accelerator & 15{,}686 instances, lower channel and left edge & this document \\
Fill, decap and antenna diodes & 107{,}890 instances, not drawn in the maps & this document \\
Cells not attributable by connectivity & 5{,}916, 9.75\,\% of the logic & this document \\
\bottomrule
\end{tabularx}
\end{table}

Two rows deserve emphasis because they are the ones a reader is most
likely to misread.

\textbf{The worst setup slack is negative and that is not a defect
report.} No frequency is claimed for this design
(\dcite{53}{}, \dcite{05}{rule 5}); the constraint the flow was given is
a period, and the slack is measured against it. A negative worst slack
at the slow corner says what period this layout would \emph{not} make
at that corner, and the honest statement of the result is the number
and the corner together, never a megahertz figure derived from either.

\textbf{The hold margin is 29.6\,ps and everything competes for it.}
That single number is the reason three separate experiments in this
project ended the same way: a change with a clearly positive effect of
its own spends more hold margin than the design has. The antenna
repair does it (\dcite{79}{}), and so does the memory read register
(\cref{sec:rdreg-result}). A design that meets hold by 29.6\,ps is a
design in which the next improvement must be paid for out of hold.

\section{Four experiments, 14--15 September 2026}

\subsection{The event engine's invariant, proved over all reachable
states}

\dcite{56}{section 9.4} had named a gap and declined to close it twice:
the largest register-transfer file in the system-on-chip carried no
formal property set, and the one-line invariant on which its event
mechanism rests --- that the strobe presenting an event to a neuron is
exactly the one state of the event-engine state machine, neither
implied by it nor implying it --- was checked by one simulation test
and one textual guard rather than over all reachable states.

It is now a property set, \code{hw/soc/formal/soc\_npu.sby}, and the
invariant holds: bounded model checking to depth 30, $k$-induction at
depth 16, and a cover showing the state is reachable, all pass.

Two details of how it was made to work are worth carrying, because they
are the kind of thing that costs a day and appears in no manual. The
properties are \emph{included into} the module under a formal guard, not
read as a separate compilation unit, so that the state machine's
signals are in scope natively; a hierarchical reference from an outside
wrapper into a generate scope silently declares a fresh free wire
instead, which is how an earlier attempt on a different block produced a
counterexample that the register-transfer code cannot produce. And the
six submodules in the instantiation closure are read \emph{without} the
formal flag, so that their own property includes --- and the
assumptions inside them --- stay out of this model.

\subsection{The register file with its storage and its scrub}

\dcite{63}{section 18} had measured that the register-file consistency
check of the RISC-V formal framework proves to check cycle 18 on the
core with the error-correcting register file substituted and times out
at 21, 23 and 25 under six engines at four hours each. That document
ranked three routes out and took none of them.

Two were taken here, and the result is a clean separation of what is
hard from what is merely large.

Route 2 asks only the register file: a wrapper that writes a word,
waits while the scrub engine walks, and reads it back, with three more
properties for the scrub's coverage, its arbitration against the core
and its silence when no fault is present. Cover passes --- the scrub
visits all 31 architectural registers. Bounded model checking and
$k$-induction do not: four engines, including a bit-level model checker
and an IC3 implementation, sit at step 4 and two of them report
\textsc{timeout} at 10{,}800 seconds. Step 4 is the first cycle at which
the read-back property examines a stored word, so the solver is being
asked to relate eight parity trees at the write encoder to eight at the
read decoder through thirty-two more encoders in the scrub path.

Route 3 removes exactly that. With the codec replaced by a degenerate
stand-in of the same interface, an IC3 engine proves all four
properties \emph{unbounded} in 44 seconds. The codec's own correctness
is a separate theorem that another job in the same directory already
proves, so the composition is sound: a property of the wrapper that
holds for every codec satisfying decode-of-encode-is-identity holds for
the real one.

The abstraction also taught the wrapper's hidden assumptions. The first
stand-in was the identity code, and bounded model checking failed at
step 3 with every read bit inverted. The cause is in the read path's
correction mask, which flips every bit whose syndrome matches its
column of the parity-check matrix: with an all-zero column set and a
zero syndrome, that is every bit. The register file therefore assumes
that every column of its parity-check matrix is non-zero and that the
encoding of the all-zero word is zero --- both true of the real code by
construction, and both now written down where the next person to
abstract this block will find them.

\subsection{The multiply and divide checks, bounded rather than
stopped}
\label{sec:mext-result}

\dcite{63}{section 8.4} reported six instruction checks stopped without
a verdict after 35 to 40 minutes each, and section 11 of the same
document listed a different solver as the first route out. That route
was taken: all six checks were run again on both the stock and the
substituted core under a second bit-blasting solver, each with an
explicit time bound written into the job so that the status the tool
records is its own word.

All twelve report \textsc{timeout} at 7{,}200 seconds. None closes.

The interesting number is not in the checks but in the witnesses. The
same six cover jobs, which the original solver reached in 9 to 73
seconds, take 1{,}388 to 10{,}999 seconds under the new one --- between
roughly twenty and one hundred and fifty times slower --- while still
passing. Whatever eventually closes a sequential multiplier
equivalence, it is not a change of satisfiability back end: two of them
have now been given the problem and the second is worse at both halves
of it.

\subsection{The memory read register, measured on the design that
exists}
\label{sec:rdreg-result}

\dcite{50}{} priced a response register on the memory read path twice,
on two designs that no longer exist, and five later documents each
recorded that the measurement was owed on the design that does. It was
taken: one synthesis with the register enabled, one full layout of that
netlist, one whole-system simulation.

The mechanism works exactly as \dcite{50}{} predicted. The macro read
return stops being a violating launch group: 360 violating endpoints
become 17, their total violation falls from $-864.88$ to $-28.87$\,ns,
and the design's total setup violation improves by 21.9\,\%, the
largest single-change improvement any layout in this project has
produced. The seventeen that remain do not reach the bus at all;
sixteen of them end in the scrub engine's counter.

And it does not make the part. The worst path is unchanged in kind and
0.16\,ns worse in value, because it was never behind a macro: it is the
register file's read address into the register file's own error-correcting
encoder, whose structural answer is not the fabric's at all. The fabric
captures 54 of the 3{,}529 violating endpoints and the time base 138,
against 1{,}224 in the register file (\code{docs/83} section 4). Worse, the 61 new
flip-flops and 683 new instances cost 325.7\,ps of hold margin at the
fast corner, and the design has 29.6\,ps. The register therefore stays
at its default of disabled, and what the experiment establishes is its
price at today's hold margin rather than a recommendation.

\subsection{A run that measured nothing, kept on purpose}

The first attempt at the experiment above completed all 67 flow steps
and produced a layout of the \emph{baseline} netlist. The synthesis
knob was set in the environment; the place-and-route script takes its
netlist from a different variable and had defaulted to the unmodified
one. Nothing in the flow objected, because the layout was a perfectly
good layout of the netlist it was given.

The run directory is kept, labelled, and its numbers are in no table.
It belongs in this chapter because it is the project's most common
failure shape in its purest form: not a wrong measurement, but a
correct measurement of the wrong thing. The countermeasure that would
have caught it in seconds --- count the mechanism in the netlist before
believing the layout --- is the same one the redundancy guards apply,
for the same reason.
